\pdfinfoomitdate=1
\pdftrailerid{}
\documentclass[11pt]{article}
\usepackage[letterpaper,margin=1in]{geometry}
\usepackage[T1]{fontenc}
\usepackage{lmodern,microtype,amsmath,amssymb,amsthm,booktabs,longtable,array,needspace}
\usepackage[hidelinks]{hyperref}
\usepackage{etoolbox}\AtBeginEnvironment{thebibliography}{\raggedright}% [write] removes the two underfull bibliography lines of the delivered build
\hypersetup{pdftitle={Report208 First maximum frequencies in ordered partitions},pdfauthor={Report208},pdfsubject={A386374 and A386375}}
\newtheorem{theorem}{Theorem}[section]
\newtheorem{lemma}[theorem]{Lemma}
\newtheorem{proposition}[theorem]{Proposition}
\newtheorem{corollary}[theorem]{Corollary}
\theoremstyle{remark}\newtheorem{remark}[theorem]{Remark}
\newcommand{\E}{\mathbb E}
\newcommand{\Prob}{\mathbb P}
\newcommand{\ceil}[1]{\left\lceil #1\right\rceil}
\newcommand{\floor}[1]{\left\lfloor #1\right\rfloor}
\newcommand{\abs}[1]{\left|#1\right|}
\newcommand{\norm}[1]{\left\|#1\right\|}
\DeclareMathOperator{\Var}{Var}
% Added in the write (batch 113, 7 October 2026).
\newcommand{\file}[1]{\nolinkurl{#1}}
\newenvironment{writenote}[1][]{\par\smallskip\begingroup\small
 \noindent\textbf{[write]}\if\relax\detokenize{#1}\relax\else~\textbf{#1.}\fi\ }%
 {\par\endgroup\smallskip}
\setlength{\emergencystretch}{2em}
\allowdisplaybreaks[2]
\title{First maximum frequencies in ordered partitions}
\author{Report208}
\date{4 October 2026}
\begin{document}
\maketitle
\begin{abstract}
A386374 counts initial-alphabet words for which letter 1 is tied for highest frequency, equivalently ordered set partitions whose first block is largest. Its strict companion is A386375. Individual positive-pole subtraction gives absolutely convergent residue sums with explicit absolute exponential-generating-function coefficient errors 11 and 4. Uniform analytic reversion then gives finite root-free expansions at every fixed order. For the weak count, a two-term phase law yields different sharp upper and lower oscillation scales, including the constant term in each phase minimum. The strict count is uniformly smaller by a factor asymptotic to $(\log 2)\log\log n/\log n$. A finite-step Newton construction and shrinking two-ceiling enclosures describe the inverse of the weak sequence. The moving-pole method and maximum-size scale are classical, with directly relevant surjection results in Gourdon's 1996 thesis. This report supplies a specialized analysis, with finite numerical validation distinguished from rigorous asymptotic statements and finite-input certification.
\end{abstract}
\begin{writenote}[Status and trust boundary, added 7 October 2026, batch~113]
This is bundle Report~208 of an external AI-assisted research session, filed
in ProveIt as the single-source research report
\file{a386374-first-block-maximum} (provenance, the sources as the write read
them, what was checked, relation to the repository, the collected non-claims
and the reading conventions are in Section~\ref{fbm:sec:provenance}). It is
unrefereed and not formalized. Its author line and PDF author field read
``Report208'': it names no person or tool. \emph{What rests on what.} Every
theorem has a conventional proof in the text: exact generating functions,
Rouch\'e's theorem and Cauchy estimates for individually subtracted poles,
analytic reversion with uniform bounds, factorial-tail estimates, Stirling's
formula and Newton's method with explicit derivative bounds. No proof uses a
computation; the exact checks corroborate finite identities and the floating
tables are illustrations. The write adds Remarks~\ref{fbm:rem:transseries}
and~\ref{fbm:rem:oeis} and dated notes; no statement, proof or number of the
source was changed.
\end{writenote}
\setcounter{tocdepth}{1}
\tableofcontents
\newpage

\section{Objects and main conclusions}
For $n\geq1$, let $a_n$ count words of length $n$ on an initial interval $[k]=\{1,\ldots,k\}$, for any $1\leq k\leq n$, using every letter and satisfying
\[
 \#\{i:w_i=1\}\geq \#\{i:w_i=j\}\qquad(1\leq j\leq k).
\]
Let $u_n$ impose strict inequality for every $j>1$, and put $a_0=u_0=1$. These are OEIS A386374 and A386375, respectively \cite{weak,strict}. The index sets of the letters form an ordered set partition. Thus the first block is tied for largest in the weak model and uniquely largest in the strict model.

Write $F_n$ for the ordered Bell number, and set
\[
 \rho=\log 2,\qquad P_n=\frac{a_n}{F_n},\qquad E_m(z)=\sum_{j=0}^m\frac{z^j}{j!},
 \qquad E_m(z)=0\quad(m<0).
\]
All logarithms are natural unless a base is displayed. Asymptotic statements concern $n\to\infty$ (or the explicitly indicated real variable) with expansion order fixed.

For a uniform random ordered partition of $[n]$, let $K$ be its number of blocks and $J$ the number of blocks tied for largest. Conditional on the unordered partition, all block orders are equally likely. Consequently
\begin{equation}\label{fbm:eq:expectations}
 P_n=\E\!\left[\frac{J}{K}\right],\qquad
 \frac{u_n}{F_n}=\E\!\left[\frac{\boldsymbol 1_{\{J=1\}}}{K}\right].
\end{equation}
In particular, $P_n\geq1/n$. The random denominator in \eqref{fbm:eq:expectations} must be retained: $\E[J/K]$ is not generally $\E[J]/\E[K]$.

The central oscillation conclusions are
\begin{equation}\label{fbm:eq:mainenvelopes}
 \limsup_{n\to\infty}\frac{nP_n\log\log n}{\log n}=\frac2e,
 \qquad
 \liminf_{n\to\infty}\frac{nP_n}{\log\log n}=2\rho.
\end{equation}
These have different normalizations. The strict/weak ratio has the phase-uniform equivalent
\begin{equation}\label{fbm:eq:mainstrict}
 \frac{u_n}{a_n}\sim\rho\frac{\log\log n}{\log n}.
\end{equation}
Sections \ref{fbm:sec:poles}--\ref{fbm:sec:expansion} give stronger finite-sum coefficient expansions; Sections \ref{fbm:sec:phase}--\ref{fbm:sec:envelopes} prove \eqref{fbm:eq:mainenvelopes}; Section \ref{fbm:sec:strict} proves \eqref{fbm:eq:mainstrict}. Section \ref{fbm:sec:inverse} concerns the inverse of $a_n$, not an inverse theorem for the strict sequence.

\subsection{Exact generating functions and recurrences}
Fix the size $m$ of the first block. The remaining labels form an ordered sequence of nonempty blocks of size at most $m$ in the weak case and at most $m-1$ in the strict case. Hence
\begin{align}
 A(z)=\sum_{n\geq0}\frac{a_nz^n}{n!}
 &=\sum_{m\geq0}\frac{z^m}{m!\,[2-E_m(z)]},\label{fbm:eq:weakEGF}\\
 U(z)=\sum_{n\geq0}\frac{u_nz^n}{n!}
 &=1+z+\sum_{m\geq1}\frac{z^{m+1}}{(m+1)!\,[2-E_m(z)]},\label{fbm:eq:strictEGF}\\
 F(z)=\sum_{n\geq0}\frac{F_nz^n}{n!}&=\frac1{2-e^z}.\label{fbm:eq:BellEGF}
\end{align}
The first two EGFs agree with the current OEIS entries. They also have the following direct exact recurrence. Let $b_{m,t}$ count ordered partitions on $t$ labels with block size at most $m$, including $b_{0,0}=1$ and $b_{0,t}=0$ for $t>0$. Then
\begin{align}
 b_{m,0}&=1,&
 b_{m,t}&=\sum_{j=1}^{\min(m,t)}\binom tj b_{m,t-j},\label{fbm:eq:boundedrec}\\
 a_n&=\sum_{m=1}^n\binom nm b_{m,n-m},&
 u_n&=\sum_{m=1}^n\binom nm b_{m-1,n-m}\quad(n\geq1).\label{fbm:eq:exactrec}
\end{align}
The empty sums use the stated conventions. The first values are shown in Table \ref{fbm:tab:exact}; all displayed values are generated from the exact recurrences and independently checked by rational EGF extraction.

% [write] inlined from the delivered tables/exact.tex (shipped as data/tables-exact.tex)
\begin{table}[htbp]\centering
\caption{Exact initial counts}\label{fbm:tab:exact}
\begin{tabular}{r r r}\toprule
$n$ & $a_n$ & $u_n$ \\\midrule
0 & 1 & 1 \\
1 & 1 & 1 \\
2 & 3 & 1 \\
3 & 10 & 4 \\
4 & 47 & 17 \\
5 & 276 & 96 \\
6 & 2022 & 652 \\
7 & 17606 & 5356 \\
8 & 179391 & 51361 \\
9 & 2093860 & 568840 \\
10 & 27581888 & 7157036 \\
11 & 404680398 & 101048454 \\
12 & 6541528886 & 1582644956 \\
\bottomrule\end{tabular}\end{table}

\subsection{Provenance, sources and reading conventions}\label{fbm:sec:provenance}
\begin{writenote}[Added 7 October 2026, batch~113]\raggedright
\emph{Provenance.} This report is Report~208 of the session bundle that
ProveIt's \file{docs/incoming} intake processed as batch~113: the archive
\file{Report208-reproducibility.zip} ($470{,}406$ bytes, 13 files, wrapper
directory \file{Report208/}), arrival commit \file{60f54ea06}, placed in
\file{a4186a946}. The text is the delivered \file{Report208.tex} (551 lines,
15 pages, dated 4~October 2026), shipped as \file{article.tex}. The programs
\file{verify.py}, \file{independent_check.py} and \file{reproduce.py} are
shipped under \file{code/}; the generated data, the requirements and the three
generated tables under \file{data/} (the tables as
\file{data/tables-exact.tex}, \file{tables-diagnostic.tex} and
\file{tables-phase.tex}); \file{SOURCES.txt} is shipped as delivered. The text
reads its tables with \verb|\input{tables/...}|; the write prints the shipped
table files inline at those places. Not shipped, and retrievable from the
arrival commit: the delivered \file{Report208.pdf}; the delivered
\file{README.txt}, which was staged as the report README and is replaced by
it. The archive carries no checksum manifest (its README says the hashes
accompany the deliverables outside the ZIP; none were delivered). The package
records no ProveIt commit; it carries no ``prepared for private review'' line
and no personal data. It is a single-source report, so no merge choice was
made. The write prefixed the 68 delivered labels (65 in the text, 3 in the
generated tables) with \file{fbm:} and updated the references to them
(44 \verb|\eqref|, 20 \verb|\ref|); it added this subsection, the notes marked [write] and
Remarks~\ref{fbm:rem:transseries} and~\ref{fbm:rem:oeis}, each the last
statement of its section, and inserted nothing before a delivered display,
statement or table, so every theorem, section, equation and table of the
source keeps its delivered number. One preamble line sets the bibliography
ragged right (\texttt{etoolbox}), which removes the two underfull lines of the
delivered build (the entries of Daly and of Carayol--Rotondo).

\emph{The sources, as the write read them} (7~October 2026): OEIS A386374,
A386375 and A308876 in the internal format, with the b-files of the first two
(Remark~\ref{fbm:rem:oeis}); the transseries volume
\path{Analysis/Transseries/docs/series-and-transseries/Transseries_And_Inversion/transseries_and_inversion.tex}
(\file{p0:def:model}, \file{p0:thm:lambert-core}, \file{p0:def:three-inverses},
\file{p0:thm:staircase}, \file{plt:thm:lw-template}). Not read by the write:
the cited literature, in particular Gourdon's thesis and proceedings paper,
Flajolet--Sedgewick, Daly, Carayol--Rotondo, Bender--Gao and Olofsson; the
source's account of them (Section~\ref{fbm:sec:sources} and
\file{SOURCES.txt}) is reported, not confirmed. The source itself records that
it did not obtain the full texts of Bender--Gao, of Gourdon's 1998 journal
version and of Olofsson.

\emph{What was checked at intake.} The batch-113 dossier read the manuscript,
found no demonstrably wrong claim, named Gourdon's 1996 thesis as the direct
prior and Bender--Gao's unread full text as an open overlap, and reran the
delivered suite on a copy (it reproduces). \emph{What the write checked.} It
read every proof and found no error. With its own code (Python 3.14.4, exact
integers and fractions; SymPy 1.14.0; mpmath 1.3.0 at 50--260 digits), not
importing the delivered programs, it
\begin{itemize}\raggedright
\item computed $a_n$ and $u_n$ from~\eqref{fbm:eq:boundedrec}--\eqref{fbm:eq:exactrec}
for $n\le425$, equal to all $426$ terms of each OEIS b-file and to
Table~\ref{fbm:tab:exact}; the weak EGF~\eqref{fbm:eq:weakEGF} by exact series
to degree~40; $a_{n+1}>a_n$ for $1\le n<425$;
\item enumerated all initial-alphabet words of length $n\le7$ ($47{,}293$ at
$n=7$): both counts and the identity $P_n=\E[J/K]$
of~\eqref{fbm:eq:expectations} hold exactly;
\item checked every numerical inequality of the proofs of
Theorems~\ref{fbm:thm:poles} and~\ref{fbm:thm:strict} and of
Section~\ref{fbm:sec:inverse} ($\sum_{j\ge5}1/j!=0.00994\ldots$,
$2e^{-1/2}\sin(1/8)=0.1512\ldots$, $(\sqrt8-1)^2/24=0.1392\ldots$,
$r_3=0.6988\ldots$, $15(e-2)=10.77\ldots$, $15(e-5/2)=3.27\ldots$,
$1/(2(1-\rho))=1.629\ldots$, the rational checks after~\eqref{fbm:eq:logerror});
on the unit circle the sampled minima of $|e^z-2|$ and $|2-E_m(z)|$,
$2\le m\le12$, are at least $1/2$, well above the bounds $1/8$ and $1/12$;
the deviations in~\eqref{fbm:eq:absweak}, \eqref{fbm:eq:absstrict}
and~\eqref{fbm:eq:Bellpole} for $n\le60$ are at most $0.37$, $0.12$
and~$0.05$, inside $11$, $4$ and~$10$;
\item rederived by its own reversion the coefficients~\eqref{fbm:eq:ellcoeff}
and~\eqref{fbm:eq:Q12} (SymPy, zero differences), and checked numerically that
$c_m(\rho/r_m)^n=T_m(n)G_{m,nD_m}(1/n)$ (relative differences below
$10^{-42}$);
\item recomputed Tables~\ref{fbm:tab:diagnostic} and~\ref{fbm:tab:phase}
(note in Section~\ref{fbm:sec:checks}), the real-index phase minimum against
the constant of~\eqref{fbm:eq:minimum}, and the strict/weak ratio of the
residue sums against~\eqref{fbm:eq:strictphase}.
\end{itemize}
It reran, on fresh copies of the shipped programs, \file{verify.py} and
\file{independent_check.py} (normal and \texttt{-O}): the outputs equal
\file{data/validation.json} and, apart from the recorded optimization flag that
the driver omits, \file{data/independent.json}; the table generator of
\file{reproduce.py}, applied to them, reproduces the three shipped table files
byte for byte. The PDF and ZIP builder was not run. Floating computations are
observations, not certificates.

\emph{Relation to the repository.} Before batch~113 no file of the repository
named A386374, A386375 or A308876. Ordered set partitions (the ordered Bell
numbers A000670) occur elsewhere in the collection with other statistics:
\file{a122399-surjection-diagonal}, \file{a261781-matrix-compositions}, and
\file{a173217-ordered-tuple-relations} (batch~113), whose Fubini poles are
the poles of~\eqref{fbm:eq:BellEGF}. None of them treats block maxima, so no
result is shared and no reciprocal note is proposed. No Lean or Rocq file
treats these sequences; no statement of this report is formalized, and its
place in the collection confers no formal status. The inverse is compared
with the transseries volume in Remark~\ref{fbm:rem:transseries}.
\end{writenote}
\begin{writenote}[Correction at the independent check, 7 October 2026]
The deviation bound $0.05$ for~\eqref{fbm:eq:Bellpole} holds for
$1\le n\le60$ (maximum $0.0406\ldots$). At $n=0$ the deviation is
$1-1/(2\rho)=0.2786\ldots$, still inside $10$; the two other bounds concern
$n\ge1$ and $n\ge2$, as their theorems do.
\end{writenote}
\begin{writenote}[What is not claimed, collected from the source]
The moving-pole method, the factorial maximum-size scale, the pole
perturbation and bounded variance of the maximum, and the shrinking transition
width are prior (Gourdon's 1996 thesis, Example~8; Flajolet--Sedgewick
Section~V.2); the constrained-Poisson representation of surjection profiles is
prior (Carayol--Rotondo); no novelty is claimed for discrete ties or Poisson
approximation of maxima (Daly, Olofsson); theorem-level non-overlap with
Bender--Gao is not established (abstract only). No worldwide novelty or
first-proof claim; the absence of an asymptotic formula from OEIS does not
establish originality. The expansion is a moving-window finite sum, not a
constant-coefficient power series in $1/n$; all orders are fixed; the
constants $C_L$ and onsets are not certified; the inverse theorem concerns the
weak sequence only; the tables use finite fixed ranges and floating
arithmetic and certify nothing; no unconditional formula
$q(X)=\lceil x_L\rceil$ is claimed. \emph{The write adds:} its checks are
exact or floating computations; its transseries remark claims no novelty for
any inversion.
\end{writenote}
\medskip
\begingroup\small
\noindent\emph{Reading conventions} (letters the source uses in more than one
sense, with the tempting false readings; no symbol was renamed).\par\smallskip
\renewcommand{\arraystretch}{1.1}
\noindent\begin{tabular}{@{}>{\raggedright\arraybackslash}p{0.8in}>{\raggedright\arraybackslash}p{5.4in}@{}}
\toprule
Symbol & Meanings\\
\midrule
$a_n$, $\alpha_m$ & $a_n$ the weak counts (A386374); $\alpha_m=E_{m-1}(\rho)$. \emph{False reading:} $\alpha_m$ is not $a_m$ (the source says so in Section~\ref{fbm:sec:expansion}).\\
$m$ & The first-block size in~\eqref{fbm:eq:weakEGF}; the pole index; the phase index $m(n)$ of Proposition~\ref{fbm:prop:simple}, with $m\sim\log n/\log\log n$.\\
$P_n$, $p_m$ & $P_n=a_n/F_n$ the tie probability; $p_m=\rho^m/m!$ the factorial weights.\\
$D$, $d$, $\Delta$ & $D_m=2-E_m(\rho)$ the factorial tails; $d_m(s)$ the reversion; $\Delta_L(N)$ the enclosure radius.\\
$B$ & $B_{m,j}(x)$ the exponent coefficients~\eqref{fbm:eq:bj}; $B(x)$ the interpolation~\eqref{fbm:eq:interpolation}; $B_L(z)$ its finite model.\\
$S$, $T$, $Q$, $q$ & $S(n)$ the normalized residue sum, $S_L$, $S_0$; $T_m(n)$ the root-free amplitudes; $Q_{m,l}$ the correction polynomials; $q(X)$ the threshold. \emph{False reading:} $q$ is not $Q$.\\
$N$, $x$, $y$, $z$ & $N$ the Lambert scale~\eqref{fbm:eq:N}; $x$ the real solution of $B(x)=X$, also the argument $x=nD_m$ of $Q_{m,l}$; $y=\log X$; $z$ the EGF variable and the real variable of Section~\ref{fbm:sec:inverse}, $z_j$ Newton iterates.\\
$t$, $v$, $\lambda$ & $t=\lambda_m(n)$ the phase coordinate; $v$ the expansion variable $1/n$ in $G_{m,x}$ and $v=\rho t/m$ in the proof of Theorem~\ref{fbm:thm:envelopes}; $\lambda_m(n)=nD_m/(2\rho)$.\\
$J$, $K$ & $J$ the number of tied largest blocks and $J_L$ the number of Newton steps; $K$ the number of blocks and $K_L=20(L+3)$ the window constant.\\
$c$, $C$ & $c_m$ the residues and $c^{\mathrm s}_m$ their strict versions; $c$, $c'$, $C$, $C_L$ constants.\\
\bottomrule
\end{tabular}
\endgroup

\section{Positive poles and an explicit absolute error}\label{fbm:sec:poles}
For $m\geq1$, let $r_m$ be the unique positive solution of $E_m(r_m)=2$, and define
\begin{equation}\label{fbm:eq:residue}
 c_m=\frac{r_m^{m-1}}{m!E_{m-1}(r_m)}.
\end{equation}
Existence and uniqueness follow from monotonicity on the positive real axis. We have $r_1=1$, $r_2=\sqrt3-1$, and
\[
 \rho<r_{m+1}<r_m\leq\sqrt3-1\quad(m\geq2),\qquad r_m\downarrow\rho.
\]
Indeed $E_m(\rho)<2$, and $E_m$ increases to $e^z$ on positive bounded intervals. Also $0<c_m\leq1/m!$, because $r_m\leq1$ and $E_{m-1}(r_m)\geq1$.

\begin{theorem}[Uniform residue estimate]\label{fbm:thm:poles}
For every integer $n\geq1$,
\begin{equation}\label{fbm:eq:absweak}
 \abs{\frac{a_n}{n!}-\sum_{m\geq1}c_mr_m^{-n}}\leq11.
\end{equation}
The series converges absolutely for every fixed real $n$. The bound is an absolute error for the EGF coefficient $a_n/n!$.
\end{theorem}

\begin{proof}
We first separate all other complex zeros uniformly on the unit circle. Write $z=x+iy$, $|z|=1$. If $|y|\geq1/4$, then
\[
 |e^z-2|\geq2e^{x/2}|\sin(y/2)|
 \geq2e^{-1/2}\sin(1/8)>\frac18.
\]
For completeness, squaring the first inequality follows from
$|e^{x+iy}-2|^2=(e^x-2)^2+8e^x\sin^2(y/2)$.
If $|y|<1/4$, then $|x|>\sqrt{15}/4$, and the reverse triangle inequality gives
$|e^z-2|\geq|e^x-2|>1/8$.
For $m\geq4$, the exponential tail on $|z|=1$ is at most
\[
 \sum_{j\geq5}\frac1{j!}<\frac1{100}.
\]
Rouch\'e's theorem now shows that $2-E_m$ has exactly one zero inside the unit disk, since $2-e^z$ has only $\rho$ there. This zero is $r_m$, and
\begin{equation}\label{fbm:eq:circle}
 |2-E_m(z)|>\frac18-\frac1{100}>\frac1{12}\qquad(|z|=1).
\end{equation}
For $m=2$, the factorization with roots $\sqrt3-1$ and $-\sqrt3-1$ gives the same lower bound directly. For $m=3$, the derivative of $z^3+3z^2+6z-6$ is strictly positive on the real axis, so it has just one real root $r_3<3/4$. The conjugate nonreal roots have modulus $\sqrt{6/r_3}>\sqrt8$. Thus the modulus of $2-E_3(z)$ on $|z|=1$ is greater than $(\sqrt8-1)^2/24>1/12$.

For each $m\geq2$, remove its own simple pole before summing:
\[
 H_m(z)=\frac{z^m}{m!\,[2-E_m(z)]}-\frac{c_m}{1-z/r_m}.
\]
This is analytic on a neighborhood of the closed unit disk. The first term on its boundary is bounded by $12/m!$. For the second, use
\[
 \frac{c_m}{|1-z/r_m|}\leq\frac1{m!}\frac{r_m}{1-r_m}
 \leq\frac{\sqrt3+1}{m!}<\frac3{m!}.
\]
It follows that $|H_m(z)|\leq15/m!$ throughout the disk. The series of these pole-subtracted functions converges normally in the disk and uniformly on its boundary. Cauchy's coefficient estimate gives
\[
 \abs{[z^n]\sum_{m\geq2}H_m(z)}\leq15(e-2)<11.
\]
For $m=1$, the summand $z/(1-z)$ and its pole term $1/(1-z)$ have the same positive-degree coefficients. The $m=0$ term affects only degree zero. This proves \eqref{fbm:eq:absweak}. Absolute convergence of the residue series follows from $c_m\leq1/m!$ and $\rho\leq r_m\leq1$ for any fixed real exponent.
\end{proof}

\begin{remark}
The positive poles accumulate at $\rho$. The argument subtracts individual poles first; it neither moves a contour through their accumulation point nor assumes meromorphic continuation across that point.
\end{remark}

\section{Conversion to a relative estimate}\label{fbm:sec:relative}
Pole subtraction in \eqref{fbm:eq:BellEGF} gives
\begin{equation}\label{fbm:eq:Bellpole}
 \frac{F_n}{n!}=\frac1{2\rho^{n+1}}+O(1).
\end{equation}
This follows directly by subtracting $1/[2\rho(1-z/\rho)]$ and using the same unit circle. In fact the boundary norm of the remainder is less than $10$: the first term is at most $8$, and the pole term at most $1/[2(1-\rho)]<5/3$, using $\rho<7/10$.

Define
\[
 S(n)=\sum_{m\geq1}c_m(\rho/r_m)^n.
\]
The lower bound $a_n/F_n\geq1/n$, \eqref{fbm:eq:absweak}, and \eqref{fbm:eq:Bellpole} imply $S(n)\geq c/n$ for some absolute $c>0$ and all sufficiently large $n$. Therefore
\begin{equation}\label{fbm:eq:relative}
 a_n=n!\rho^{-n}S(n)\,[1+O(n\rho^n)].
\end{equation}
The original constant 11 is not an $O(1)$ error in $a_n$, and the conversion to a relative error requires the lower bound just proved.

\section{Finite root free expansions at every fixed order}\label{fbm:sec:expansion}
Set
\begin{equation}\label{fbm:eq:D}
 D_m=2-E_m(\rho)=\sum_{j>m}\frac{\rho^j}{j!},\quad
 \alpha_m=E_{m-1}(\rho),\quad \beta_m=E_{m-2}(\rho),\quad
 \gamma_m=E_{m-3}(\rho).
\end{equation}
The Greek notation avoids confusing the truncated-exponential value with the count $a_m$. Let $d_m(s)$ denote the local analytic solution of
\[
 E_m(\rho+d_m(s))-E_m(\rho)=s,\qquad d_m(0)=0.
\]
Define coefficient series
\begin{align}
 \ell_m(s)&=\log(1+d_m(s)/\rho)=\sum_{j\geq1}\ell_{m,j}s^j,\label{fbm:eq:ell}\\
 h_m(s)&=(m-1)\ell_m(s)-\log\frac{E_{m-1}(\rho+d_m(s))}{\alpha_m}
       =\sum_{j\geq1}h_{m,j}s^j.\label{fbm:eq:h}
\end{align}
All coefficient operations at any fixed order are finite. In particular $\ell_{m,1}=1/(\rho\alpha_m)$. For $x\geq0$, put
\begin{align}
 B_{m,j}(x)&=h_{m,j}x^j-\ell_{m,j+1}x^{j+1},\label{fbm:eq:bj}\\
 Q_{m,0}(x)&=1,\qquad
 Q_{m,l}(x)=\frac1l\sum_{j=1}^l jB_{m,j}(x)Q_{m,l-j}(x)\quad(l\geq1).\label{fbm:eq:Qrec}
\end{align}
This is the coefficient recurrence for the exponential of $\sum_{j\geq1}B_{m,j}(x)v^j$.

For an integer $L\geq0$, choose
\begin{align}
 K_L&=20(L+3),\notag\\
 M_L(n)&=\min\{M\geq2:\rho^M/M!\leq n^{-L-5}\},\label{fbm:eq:cutoffs}\\
 \mathcal W_L(n)&=\{2\leq m\leq M_L(n):nD_m\leq K_L\log n\},\notag\\
 T_m(n)&=\frac{\rho^{m-1}}{m!\alpha_m}
             \exp\!\left(-\frac{nD_m}{\rho\alpha_m}\right).\label{fbm:eq:T}
\end{align}
The window has $O_L(\log n/\log\log n)$ indices.

\begin{theorem}[Uniform finite expansion]\label{fbm:thm:expansion}
For every fixed $L\geq0$,
\begin{equation}\label{fbm:eq:allorders}
 a_n=n!\rho^{-n}
 \sum_{m\in\mathcal W_L(n)}T_m(n)
 \left(\sum_{l=0}^L\frac{Q_{m,l}(nD_m)}{n^l}\right)
 \left[1+O_L\!\left(\frac{(\log n)^{2L+2}}{n^{L+1}}\right)\right].
\end{equation}
The final bracket multiplies the entire finite sum. For sufficiently large $n$ the zeroth-order sum is positive and the correction in parentheses is uniformly $1+o(1)$. No root $r_m$ needs to be calculated in this formula.
\end{theorem}

\begin{proof}
The proof has a local uniform estimate and two discarded tails. On $|d|\leq1/10$, every derivative of $E_m(\rho+d)$ is bounded in modulus by $e^{\rho+1/10}<3$, while $\alpha_m\geq1$. On a sufficiently small fixed $s$-disk, Rouch\'e's theorem applied to $\alpha_md-s$ and the quadratic remainder gives a unique analytic inverse with $|d_m(s)|\leq2|s|$, uniformly in $m$. On a smaller fixed disk, $E_{m-1}(\rho+d_m(s))$ is uniformly separated from zero. Therefore
\begin{equation}\label{fbm:eq:localbounds}
 \abs{\ell_m(s)-\frac{s}{\rho\alpha_m}}\leq C|s|^2,
 \qquad |h_m(s)|\leq Cm|s|.
\end{equation}
For fixed $m$ and $x\geq0$, form
\[
 G_{m,x}(v)=\exp\!\left(h_m(xv)-
       \frac{\ell_m(xv)-\ell_{m,1}xv}{v}\right),
\]
with its removable value at $v=0$. On $|v|=c/(m+x)^2$, this function is uniformly bounded for an absolute sufficiently small $c>0$. Cauchy's estimate at $v=1/n$ gives
\begin{equation}\label{fbm:eq:Gbound}
 G_{m,x}(1/n)=\sum_{l=0}^L Q_{m,l}(x)n^{-l}
       +O_L\!\left(\left[\frac{(m+x)^2}{n}\right]^{L+1}\right).
\end{equation}
At $x=nD_m$, the exact normalized residue is $T_m(n)G_{m,nD_m}(1/n)$. In the retained window, $m+nD_m=O_L(\log n)$. Summing \eqref{fbm:eq:Gbound} against the positive $T_m(n)$ transfers its error to a relative error for the retained zeroth-order sum.

For the left tail, the mean value theorem on $[\rho,1]$ gives $r_m-\rho\geq D_m/e$, while
\[
 \log(r_m/\rho)\geq(r_m-\rho)/r_m\geq D_m/e.
\]
Hence indices with $nD_m>K_L\log n$, including $m=1$ for large $n$, contribute at most $e n^{-K_L/e}$ to $S(n)$. This is smaller than the needed $n^{-L-3}$. The corresponding zeroth-order root-free tail is no larger in order, since $\rho\alpha_m\leq2\rho<e$.

For the right tail, $r_m-\rho\leq D_m$ and
$(r_m/\rho)^{m-1}\leq\exp((m-1)D_m/\rho)\leq C$ uniformly. Thus $c_m\leq C\rho^{m-1}/m!$, and the tail past $M_L(n)$ is $O(n^{-L-5})$. After division by $S(n)\geq c/n$, both tails remain smaller than the claimed relative error. Finally, \eqref{fbm:eq:relative} adds an exponentially small term. The $L=0$ case of the uniform local bound also shows that the retained corrected sum is positive and comparable to its zeroth-order version, establishing the multiplicative form in \eqref{fbm:eq:allorders}.
\end{proof}

\subsection{The first two correction polynomials}
Writing $\alpha=\alpha_m$, $\beta=\beta_m$, $\gamma=\gamma_m$, the required coefficients are
\begin{align}
 \ell_1&=\frac1{\rho\alpha},&
 \ell_2&=-\frac{\beta}{2\rho\alpha^3}-\frac1{2\rho^2\alpha^2},\notag\\
 \ell_3&=\frac{3\beta^2-\alpha\gamma}{6\rho\alpha^5}
       +\frac{\beta}{2\rho^2\alpha^4}+\frac1{3\rho^3\alpha^3},\label{fbm:eq:ellcoeff}\\
 h_1&=\frac{m-1}{\rho\alpha}-\frac{\beta}{\alpha^2},&
 h_2&=(m-1)\ell_2-\frac{\gamma}{2\alpha^3}+\frac{\beta^2}{\alpha^4}.\notag
\end{align}
Consequently
\begin{equation}\label{fbm:eq:Q12}
 Q_{m,1}(x)=h_1x-\ell_2x^2,\qquad
 Q_{m,2}(x)=h_2x^2-\ell_3x^3+\tfrac12Q_{m,1}(x)^2.
\end{equation}
The reproduction code checks formal reversion and these coefficients symbolically, and checks recurrence \eqref{fbm:eq:Qrec} against exponential expansion through order four. This moving-window expansion provides arbitrary fixed algebraic relative precision, with its logarithmic factors explicit; it is not a constant-coefficient power series in $1/n$.

\section{A simpler sum and a two term phase law}\label{fbm:sec:phase}
Put
\[
 p_m=\frac{\rho^m}{m!},\qquad \lambda_m(n)=\frac{nD_m}{2\rho}.
\]
\begin{proposition}\label{fbm:prop:simple}
The weak probability satisfies
\begin{equation}\label{fbm:eq:simple}
 P_n=\left(\sum_{m\geq1}p_me^{-\lambda_m(n)}\right)
       \left[1+O\!\left(\frac{(\log n)^3}{n}\right)\right].
\end{equation}
For sufficiently large $n$, let $m=m(n)$ be the unique integer with
$\lambda_m\geq1>\lambda_{m+1}$, and set $t=\lambda_m$. Uniformly over this full phase,
\begin{equation}\label{fbm:eq:twoterm}
 P_n=(p_me^{-t}+p_{m+1}e^{-\lambda_{m+1}})\,[1+O(1/m)],
 \qquad m\sim\frac{\log n}{\log\log n}.
\end{equation}
\end{proposition}

\begin{proof}
Apply Theorem \ref{fbm:thm:expansion} with $L=0$. On its effective window,
$2-\alpha_m=D_{m-1}=D_m+p_m=O(mD_m)$. Replacing the residue amplitude by the corresponding factor $p_m$ after normalization by $F_n$ changes it relatively by $O(mD_m)$. Replacing the exponential rate $1/(\rho\alpha_m)$ by $1/(2\rho)$ changes the exponent by $O(nmD_m^2)$. Both are $O((\log n)^3/n)$ on the window. The discarded tails have already been bounded, and \eqref{fbm:eq:Bellpole} contributes an exponentially small relative error. This proves \eqref{fbm:eq:simple}.

The factorial tail obeys $D_m\sim p_{m+1}$, so the two inequalities defining $m$ and Stirling's formula give $m\log m\sim\log n$.
For the upper-index tail,
\[
 \sum_{j\geq m+2}p_je^{-\lambda_j}\leq\sum_{j\geq m+2}p_j
       =O(p_{m+1}/m),
 \qquad p_{m+1}e^{-\lambda_{m+1}}\geq p_{m+1}/e.
\]
For the lower-index tail, split $j<m/2$ and $m/2\leq j<m$. In the first range, $D_j\geq D_{\floor{m/2}}$, and the factor $e^{-nD_j/(2\rho)}$ is superexponentially small in $m$ relative to the retained terms. In the second range,
\[
 \frac{D_{j-1}}{D_j}\geq cm,\qquad \frac{p_{j-1}}{p_j}\leq Cm.
\]
Starting from $\lambda_m=t\geq1$, successive downward summands have ratio at most $Cm e^{-cm}$, so their sum is negligible. These bounds prove the uniform relative error in \eqref{fbm:eq:twoterm}, including the ends of each phase.
\end{proof}

\section{Sharp upper and lower oscillation scales}\label{fbm:sec:envelopes}
The elementary factorial-tail expansions are
\begin{align}
 D_m&=p_{m+1}\left[1+\frac{\rho}{m+2}+O(m^{-2})\right],\label{fbm:eq:tailratio}\\
 \frac{D_{m+1}}{D_m}&=\frac{\rho}{m+2}[1+O(m^{-2})],\notag\\
 np_m&=2(m+1)t[1+O(1/m)],\qquad
 np_{m+1}=2\rho t[1+O(1/m)].\notag
\end{align}
They hold uniformly as $m\to\infty$.

\begin{theorem}[Envelopes and refined phase minima]\label{fbm:thm:envelopes}
The limits in \eqref{fbm:eq:mainenvelopes} hold. More precisely, over each full integer phase interval,
\begin{equation}\label{fbm:eq:minimum}
 \min_{\substack{n\in\mathbb Z\\2\rho/D_m\leq n<2\rho/D_{m+1}}}nP_n
 =2\rho\,[\log m+\log\log m+1-\log\rho+o(1)].
\end{equation}
\end{theorem}

\begin{proof}
For the upper envelope divide \eqref{fbm:eq:twoterm}, multiplied by $n$, by $m$. Its terms are asymptotically
\[
 2te^{-t}\quad\hbox{and}\quad 2ve^{-v},\qquad
 v=\frac{\rho t}{m}\in[0,1+o(1)].
\]
If $t$ stays bounded, the second tends to zero and the first is at most $2/e$. If $t\to\infty$, the first tends to zero and the second is at most $2/e$. Subsequence splitting covers arbitrary choices of phases. To attain the bound choose
\begin{equation}\label{fbm:eq:peak}
 n_m=\ceil{\frac{2\rho}{D_m}}.
\end{equation}
Then $n_m$ is in phase $m$ for large $m$ and $t\to1$. The ceiling is important for this literal phase statement: nearest-integer rounding could fall just below the boundary. Since $m\sim\log n/\log\log n$, the upper limit in \eqref{fbm:eq:mainenvelopes} follows.

For the lower envelope fix $\varepsilon>0$. If $t\leq(1-\varepsilon)\log m$, the first term $np_me^{-t}$ is much larger than $\log m$. If $t\geq(1-\varepsilon)\log m$ and $t=o(m)$, the second term is at least $2\rho(1-\varepsilon+o(1))\log m$. If $t$ is not $o(m)$, a subsequence has its second term of order $m$, again much larger than $\log m$. Letting $\varepsilon\downarrow0$ proves the required lower bound. For equality at leading scale, take an integer nearest to
\begin{equation}\label{fbm:eq:trough}
 \frac{2\rho}{D_m}\,[\log m+\log\log m-\log\rho].
\end{equation}
This is well inside phase $m$ for large $m$. The first term of $nP_n$ is asymptotic to $2\rho$, and its second term is
$2\rho[\log m+\log\log m-\log\rho+o(1)]$. Thus the lower limit follows from $\log m\sim\log\log n$.

For the constant term in \eqref{fbm:eq:minimum}, \eqref{fbm:eq:trough} first gives an upper bound $2\rho[\log m+\log\log m+O(1)]$. A minimizer cannot have $t$ comparable to $m$, and the second term then forces $t=O(\log m)$. More precisely it bounds $t$ above by $\log m+O(\log\log m)$; the first term bounds $t$ below by $\log m-O(1)$. On this localized range, the relative $O(1/m)$ error in \eqref{fbm:eq:twoterm} becomes absolute $o(1)$ in $nP_n$, and the upper tail is $O(\log m/m)$. Uniformly the objective reduces to
\begin{equation}\label{fbm:eq:localmin}
 2[mt e^{-t}+\rho t]+o(1).
\end{equation}
Its relevant stationary point is the large solution of
$(t-1)e^{-t}=\rho/m$; the other solution near 1 is a local maximum. The large solution has
\[
 t=\log m+\log\log m-\log\rho+o(1),\qquad
 mt e^{-t}=\rho\frac{t}{t-1}=\rho+o(1).
\]
The phase endpoints have much larger values, while the mesh in $t$ is $D_m/(2\rho)=o(1)$. This proves the claimed integer minimum with the $+1$ constant.
\end{proof}

\section{The strict first maximum companion}\label{fbm:sec:strict}
The strict EGF \eqref{fbm:eq:strictEGF} has the same denominators as the weak EGF. Its pole amplitude for index $m$ is
\[
 c_m^{\mathrm s}=\frac{r_m}{m+1}c_m.
\]
\begin{theorem}[Strict coefficient and relative estimates]\label{fbm:thm:strict}
For every integer $n\geq2$,
\begin{equation}\label{fbm:eq:absstrict}
 \abs{\frac{u_n}{n!}-\sum_{m\geq1}\frac{r_m}{m+1}c_mr_m^{-n}}\leq4.
\end{equation}
Uniformly over the phase $m=m(n)$ of Proposition \ref{fbm:prop:simple},
\begin{equation}\label{fbm:eq:strictphase}
 \frac{u_n}{a_n}=\frac{\rho}{m}[1+O(1/m)].
\end{equation}
At every fixed order, Theorem \ref{fbm:thm:expansion} extends to $u_n$ by replacing $T_m(n)$ with $\rho T_m(n)/(m+1)$ and $h_m(s)$ with $h_m(s)+\ell_m(s)$ throughout the definitions of $B_{m,j}$ and $Q_{m,l}$. The window and relative error order are unchanged.
\end{theorem}

\begin{proof}
For $m\geq2$, pole subtraction in the strict summand has unit-circle norm at most $15/(m+1)!$, exactly as in Theorem \ref{fbm:thm:poles}. Summing gives $15(e-5/2)<4$. The $m=1$ term is $z^2/[2(1-z)]$, whose coefficients equal its pole contribution $1/2$ for $n\geq2$; the polynomial $1+z$ changes only $n=0,1$. This proves \eqref{fbm:eq:absstrict}.

On the two effective indices $j=m,m+1$,
\[
 \frac{r_j}{j+1}=\frac{\rho}{m}[1+O(1/m)],
\]
because $r_j-\rho\leq D_j$ decays faster than every inverse power of $m$. The weak lower-index tail is $O(\operatorname{poly}(m)e^{-cm})$ relative to the retained terms, and remains negligible after the additional factor $m$ needed to compare with the smaller strict count. This is justified by the same $m/2$ split as in Proposition \ref{fbm:prop:simple}. On $j\geq m+2$ the strict multiplier is $O(1/m)$, so the strict upper tail is $O(1/m)$ relative to the strict two-term sum itself. Thus the strict normalized residue sum is at least $c'/(nm)$, and the absolute errors 11 and 4 become relative errors $O(nm\rho^n)$, negligible uniformly. This proves \eqref{fbm:eq:strictphase} and \eqref{fbm:eq:mainstrict}.

Finally, the factor $r_m/(m+1)$ is exactly $\rho e^{\ell_m(D_m)}/(m+1)$, which proves the announced amplitude and $h$ substitutions. Uniform analytic bounds persist. The strict lower bound $c'/(nm)$ and the generous $n^{-L-5}$ cutoff absorb the omitted factorial and left exponential tails. Hence the all-fixed-order relative error is unchanged.
\end{proof}

\section{Real inverse and shrinking integer enclosures}\label{fbm:sec:inverse}
Define the positive interpolation
\begin{equation}\label{fbm:eq:interpolation}
 B(x)=\Gamma(x+1)\sum_{m\geq1}c_mr_m^{-x},\qquad x\geq1.
\end{equation}
The sum and every fixed-order derivative converge locally uniformly. Each summand has logarithmic derivative $\psi(x+1)-\log r_m>0$ on $x\geq1$, so $B$ is strictly increasing and tends to infinity. Here $\psi$ is the logarithmic derivative of $\Gamma$. Theorem \ref{fbm:thm:poles} says $|a_n-B(n)|\leq11n!$; the interpolation is not an exact identity at integer arguments.

The sequence $a_n$ is strictly increasing for $n\geq1$. Indeed appending letter 1 is an injection into the next counted class, and $(1,2,\ldots,n+1)$ lies outside its image. The only initial equality is $a_0=a_1=1$. Thus for real $X>0$ define
\[
 q(X)=\min\{n\geq0:a_n\geq X\};\qquad q(X)=0\quad(0<X\leq1).
\]
The recurrence \eqref{fbm:eq:exactrec} settles any finite prefix outside the asymptotic range below.

\subsection{Lambert scale and a finite approximation}
Let $y=\log X$ and, for sufficiently large $X$, define
\begin{equation}\label{fbm:eq:N}
 N=\frac{y}{W(y/(e\rho))},\qquad N\log\frac{N}{e\rho}=y,
\end{equation}
where $W$ is the positive real Lambert function. The positive-sum phase estimates extend unchanged to real arguments, and Stirling's formula gives
\begin{equation}\label{fbm:eq:logB}
 \log B(x)=x\log\frac{x}{e\rho}-\frac12\log x+O(\log\log x).
\end{equation}
For example, the phase bounds give $x^{-1}$ times factors between constant multiples of $\log\log x$ and $\log x/\log\log x$ for the normalized residue sum. Their logarithms are $-\log x+O(\log\log x)$, which combined with the $\tfrac12\log x$ Stirling term yields \eqref{fbm:eq:logB}. Consequently the unique solution $B(x)=X$ satisfies
\begin{equation}\label{fbm:eq:xlocation}
 x=N+\frac12+O\!\left(\frac{\log\log N}{\log N}\right).
\end{equation}

Freeze the finite window at $\mathcal W_L(N)$ and put, for real $z$ near $N$,
\begin{align}
 R_{m,L}(z)&=\sum_{l=0}^L\frac{Q_{m,l}(zD_m)}{z^l},\notag\\
 S_L(z)&=\sum_{m\in\mathcal W_L(N)}
 \frac{\rho^{m-1}}{m!\alpha_m}e^{-zD_m/(\rho\alpha_m)}R_{m,L}(z),\label{fbm:eq:finiteinverse}\\
 B_L(z)&=\Gamma(z+1)\rho^{-z}S_L(z),\qquad g_L(z)=\log B_L(z).\notag
\end{align}
The proof of Theorem \ref{fbm:thm:expansion}, including both tails, is uniform on $I_N=[N-2,N+3]$ with this frozen window because $z/N=1+O(1/N)$. Hence
\begin{equation}\label{fbm:eq:relativeinverse}
 \frac{B_L(z)}{B(z)}=1+O_L(\epsilon_L),\qquad
 \epsilon_L=\frac{(\log N)^{2L+2}}{N^{L+1}}.
\end{equation}

\subsection{Direct derivative control}
It is unnecessary to differentiate a relative asymptotic remainder in \eqref{fbm:eq:relativeinverse}. Instead differentiate the explicit finite polynomial sum. Uniform analyticity from Section \ref{fbm:sec:expansion} gives $|h_{m,j}|\leq C_jm$ and $|\ell_{m,j}|\leq C_j$. Taking absolute values of all polynomial coefficients, recurrence \eqref{fbm:eq:Qrec} yields
\[
 |Q_{m,l}|_{\rm abs}(X)\leq C_l(m+X)^{2l},\qquad X\geq0.
\]
Each monomial of $Q_{m,l}(X)$ has degree between $l$ and $2l$. Thus $Q_{m,l}(zD_m)/z^l$ is a polynomial in $z$ of degrees between 0 and $l$. On $I_N$, its absolute monomial sum is
$O_L((\log^2N/N)^l)$, and its first and second derivatives are bounded by $l/z$ and $l(l-1)/z^2$ times that sum. It follows that
\begin{align}
 R_{m,L}&=1+O_L(\log^2N/N),\notag\\
 R'_{m,L}&=O_L(\log^2N/N^2),\qquad
 R''_{m,L}=O_L(\log^2N/N^3).\label{fbm:eq:Rderivs}
\end{align}
Let $S_0$ denote the retained positive zeroth-order sum. Since $D_m/(\rho\alpha_m)=O_L(\log N/N)$ in the window,
\[
 \frac{S_L}{S_0}=1+o(1),\quad
 \frac{|S_L'|}{S_0}=O_L(\log N/N),\quad
 \frac{|S_L''|}{S_0}=O_L(\log^2N/N^2).
\]
Together with the standard real Stirling derivative estimates these give
\begin{equation}\label{fbm:eq:gderivs}
 g_L'(z)=\log(N/\rho)+O_L(\log N/N)>0,\qquad
 g_L''(z)=O_L(1/N).
\end{equation}
The exact derivative $g'=(\log B)'$ has the same first estimate: it is $\psi(z+1)-\log\rho$ minus the positive pole-weighted mean of $\log(r_m/\rho)$. Splitting at $zD_m=K\log z$, with fixed sufficiently large $K$, bounds this mean by $O(\log z/z)$, using the left-tail estimate and the lower bound for the residue sum.

\begin{theorem}[Finite-step real inverse]\label{fbm:thm:inverse}
For fixed $L\geq0$ and all sufficiently large $X$, $B_L$ has a unique crossing $x_L$ of $X$ in $I_N$, and
\begin{equation}\label{fbm:eq:inverseerror}
 |x_L-x|\leq C_L\frac{(\log N)^{2L+1}}{N^{L+1}}
\end{equation}
for some constant $C_L$. Start $z_0=N+1/2$ and perform
\begin{equation}\label{fbm:eq:Newton}
 z_{j+1}=z_j-\frac{g_L(z_j)-\log X}{g_L'(z_j)},\qquad
 j=0,\ldots,J_L-1,\qquad J_L=\ceil{\log_2(L+2)}.
\end{equation}
After enlarging $C_L$, the same bound holds with $z_{J_L}$ in place of $x_L$.
\end{theorem}

\begin{proof}
Uniform closeness \eqref{fbm:eq:relativeinverse}, the crossing location \eqref{fbm:eq:xlocation}, and positive slope \eqref{fbm:eq:gderivs} imply existence and uniqueness in $I_N$ and the mean-value bound \eqref{fbm:eq:inverseerror}. Initially $|z_0-x_L|=O(\log\log N/\log N)$. The direct derivative bounds give the Newton error recurrence
\[
 |z_{j+1}-x_L|\leq\frac{C_L}{N\log N}|z_j-x_L|^2.
\]
For fixed $j$ all iterates stay in $I_N$ for sufficiently large $N$, and induction gives
\begin{equation}\label{fbm:eq:Newtonerror}
 |z_j-x_L|=O_{L,j}\!\left(
 \frac{(\log\log N)^{2^j}}
 {N^{2^j-1}(\log N)^{2^{j+1}-1}}\right).
\end{equation}
The chosen $J_L$ has $2^{J_L}-1\geq L+1$, so this is no larger than the order in \eqref{fbm:eq:inverseerror}. The theorem specifies a fixed finite number of steps, not unrestricted Newton convergence from arbitrary starts.
\end{proof}

\subsection{The additional integer error}
For $n\geq30$, the explicit Bell remainder bound in Section \ref{fbm:sec:relative} and $P_n\geq1/n$ imply
\begin{align}
 \frac{B(n)}{n!}&\geq\frac{\rho^{-n}}{2\rho n}-\frac{10}{n}-11
 \geq\frac{\rho^{-n}}{4\rho n},\label{fbm:eq:explicitB}\\
 \abs{\frac{a_n}{B(n)}-1}&\leq44\rho n\rho^n<31n\rho^n<\frac12,\notag\\
 |\log a_n-\log B(n)|&\leq62n\rho^n.\label{fbm:eq:logerror}
\end{align}
For the last inequality in \eqref{fbm:eq:explicitB}, one may use $\rho<7/10$ and the exact rational base-case check
\[
 (10/7)^{30}>(14/5)(10+11\cdot30).
\]
Thereafter the exponential-to-linear ratio increases. Likewise $31\cdot30(7/10)^{30}<1/2$, and $n(7/10)^n$ decreases in this range. Finally $|\log(1+u)|\leq2|u|$ for $|u|<1/2$ proves \eqref{fbm:eq:logerror}.

For large $N$, the threshold and its neighboring integers lie in $I_N$. This follows by comparing integers just below $N-1$ and just above $N+2$ with \eqref{fbm:eq:xlocation} and \eqref{fbm:eq:logerror}, and using strict monotonicity. Moreover $g'\geq\tfrac12\log N$ on $I_N$. Define
\begin{equation}\label{fbm:eq:Delta}
 \Delta_L(N)=C_L\frac{(\log N)^{2L+1}}{N^{L+1}}
       +\frac{124(N+3)\rho^{N-2}}{\log N}.
\end{equation}
The second term comes from translating \eqref{fbm:eq:logerror} to a real-index displacement using this slope bound.

\begin{corollary}[Two-ceiling enclosure]\label{fbm:cor:ceilings}
With $C_L$ sufficiently large and for all sufficiently large $X$,
\begin{equation}\label{fbm:eq:ceilings}
 \ceil{x_L-\Delta_L(N)}\leq q(X)\leq\ceil{x_L+\Delta_L(N)}.
\end{equation}
The same conclusion holds with the specified Newton iterate $z_{J_L}$ after enlarging $C_L$ to absorb \eqref{fbm:eq:Newtonerror}.
\end{corollary}

\begin{remark}[Certification boundary]
The constants $C_L$ and their validity thresholds have not been numerically certified here. Thus \eqref{fbm:eq:ceilings} is a rigorous asymptotic enclosure theorem, not automatically a certified answer at a supplied finite $X$. Since $\Delta_L(N)\to0$, its two ceilings eventually agree or differ by one. An unconditional formula $q(X)=\ceil{x_L}$ would be unjustified for $X$ arbitrarily close to an integer crossing. Certified separation or exact neighboring values from \eqref{fbm:eq:exactrec} resolve that ambiguity. Alternatively, \eqref{fbm:eq:absweak} gives explicit direct envelopes $B(n)\pm11n!$; using them as a numerical certificate requires rigorous evaluation of the residue sum, not merely floating-point roots.
\end{remark}

\begin{remark}[{[write]} The inverse and the transseries volume, 7~October 2026]\label{fbm:rem:transseries}
Compared statement by statement with the repository's transseries
volume~\cite{TSvol}; the volume's symbols are not used here.
\begin{enumerate}\raggedright
\item \emph{The growth.} $\log a_n=n\log n-(1+\log\rho)n-\frac12\log n+O(\log\log n)$
contains the term $n\log n$ and the logarithm of $nP_n$, which oscillates
between the scales $\log\log n$ and $\log n/\log\log n$
(Sections~\ref{fbm:sec:phase}--\ref{fbm:sec:envelopes}),
so it is not of the exponential--power form of \file{p0:def:model}; no
instance of the volume's model is claimed.
\item \emph{The Lambert scale.} With $w=\log(N/(e\rho))$, the
equation $N\log(N/(e\rho))=y$ of~\eqref{fbm:eq:N} becomes, after one more
logarithm, $w+\log w=\log(y/(e\rho))$: the volume's dominant block
\file{p0:thm:lambert-core} with unit slope, unit logarithmic coefficient and
target $\log(y/(e\rho))$, whose positive-coefficient branch gives exactly
$w=W(y/(e\rho))$, that is the displayed $N$. An exact instance (an equation
in $N$ for a given level).
\item \emph{The real inverse and the Newton iterates.} The location
$x=N+\frac12+O(\log\log N/\log N)$ of~\eqref{fbm:eq:xlocation} and the
iterates of Theorem~\ref{fbm:thm:inverse} rest on $\log S$, which is known
only to $O(\log\log x)$ and carries the phase oscillation, and on a finite sum
over a moving window; in the chart of the logarithm of the unknown these terms
are not a series in its reciprocal with coefficients polynomial in its
logarithm. So the write has not shown them to be instances of
\file{plt:thm:lw-template}; they are proved directly.
\item \emph{The thresholds.} $q(X)$ is the volume's staircase
(\file{p0:def:three-inverses}(1)) of the sequence $a_n$, strictly increasing
from $n=1$. The interpolation $B$ is continuous and strictly increasing but
does not take the values $a_n$ at the integers (it differs by up to $11\,n!$),
so it is not an admissible interpolation in the volume's sense, and
Corollary~\ref{fbm:cor:ceilings} is proved directly: an analogue of
\file{p0:thm:staircase}(2), not an instance, with the caution of that theorem
stated in the Certification boundary remark.
\end{enumerate}
\end{remark}

\section{Reproducible finite checks and numerical diagnostics}\label{fbm:sec:checks}
The archive contains the editable \LaTeX\ source, this PDF, two finite-check programs, a reproduction driver, dependency versions, data, and generated tables. It contains no third-party books or papers. The checks use explicit exceptions rather than Python \texttt{assert}, so optimization does not remove them.

The first program verifies all 22 displayed weak OEIS values and all 23 displayed strict values, compares composition counts through $n=12$, checks the weak counts by unordered set-partition enumeration through $n=9$, and checks formal reversion, amplitudes, and the first two correction polynomials. The second independently enumerates all initial-alphabet words through $n=7$, extracts rational EGF coefficients through $n=30$, verifies the circle and inverse constants with rational arithmetic, and checks the correction recurrence, polynomial degree ranges, and strict amplitude shift through order four. The reproduction driver compares the independent exact terms, generates every table from the data, and tests that a deliberately false guard raises an exception in both ordinary and optimized execution.

Table \ref{fbm:tab:diagnostic} is a floating-point diagnostic at 100 decimal digits. Its reference is a truncated normalized positive-pole sum over $2\leq m\leq60$; its approximations sum the $L=0,1,2$ root-free terms over the same range. This deliberately fixed diagnostic range is not the moving window of Theorem \ref{fbm:thm:expansion}. The omitted $m=1$ contribution is exponentially small for these normalized large-$n$ comparisons. The accompanying data also record small-$n$ absolute coefficient residuals after adding back the exact $m=1$ contribution. Neither diagnostic proves the error bounds or provides interval-certified roots.

% [write] inlined from the delivered tables/diagnostic.tex (shipped as data/tables-diagnostic.tex)
\begin{table}[htbp]\centering\small
\caption{Relative errors of the fixed-range root-free diagnostics}\label{fbm:tab:diagnostic}
\begin{tabular}{r r r r}\toprule
$n$ & Order $0$ & Order $1$ & Order $2$ \\\midrule
20 & $-0.0489$ & $0.000525$ & $6.04\times10^{-6}$\\
40 & $-0.05$ & $-0.000452$ & $3.75\times10^{-5}$\\
80 & $-0.0268$ & $-0.000838$ & $5.71\times10^{-6}$\\
120 & $-0.0161$ & $-0.000328$ & $-8.44\times10^{-6}$\\
1000 & $-0.00281$ & $-7.05\times10^{-6}$ & $-5.62\times10^{-8}$\\
10000 & $-0.000402$ & $-7.9\times10^{-8}$ & $-5.63\times10^{-11}$\\
1000000 & $-5.24\times10^{-6}$ & $-1.33\times10^{-11}$ & $-8.36\times10^{-17}$\\
\bottomrule\end{tabular}\end{table}

Table \ref{fbm:tab:phase} illustrates the envelope shapes. Here $n$ is the positive real argument $2\rho t/D_m$, not integer-rounded, with $t=1$ for the displayed peak and $t=\log m+\log\log m-\log\rho$ for the displayed trough. The proxy uses the normalized residue sum over $2\leq j\leq m+30$, multiplied by $2\rho$. Thus the finite rows need not equal the eventual envelope constants $2/e\approx0.7357588823$ and $2\rho\approx1.3862943611$. In particular, the slow logarithmic corrections in \eqref{fbm:eq:minimum} are visible in the trough column.

% [write] inlined from the delivered tables/phase.tex (shipped as data/tables-phase.tex)
\begin{table}[htbp]\centering
\caption{Real-index phase diagnostics}\label{fbm:tab:phase}
\begin{tabular}{r r r r}\toprule
$m$ & $\log_{10}n$ at peak & Peak $nP/m$ & Trough $nP/\log m$ \\\midrule
10 & 9.468215 & 0.894111912 & 2.694876982 \\
20 & 23.178981 & 0.815474077 & 2.502266294 \\
40 & 56.185216 & 0.775811049 & 2.397371990 \\
80 & 133.794520 & 0.755843447 & 2.320803681 \\
\bottomrule\end{tabular}\end{table}
\begin{writenote}[The tables recomputed, 7 October 2026]
The write recomputed Tables~\ref{fbm:tab:diagnostic} and~\ref{fbm:tab:phase}
with its own implementation (roots $r_m$ by Newton's method at 100 and 260
digits; root-free terms from its own coefficients $\ell_{m,j}$, $h_{m,j}$) and
found every printed value to be the correct rounding of the recomputed one,
which also equals \file{data/validation.json} to all its digits. The printed
values are roundings; the truncations that differ, at the printed number of
digits, are: in Table~\ref{fbm:tab:diagnostic}, $0.000524$ and
$6.03\times10^{-6}$ ($n=20$), $-0.0499$ ($n=40$, printed $-0.05$), $-0.0160$,
$-0.000327$ and $-8.43\times10^{-6}$ ($n=120$), $-0.00280$ and
$-5.61\times10^{-8}$ ($n=1000$), $-5.23\times10^{-6}$ and
$-1.32\times10^{-11}$ ($n=10^6$); in Table~\ref{fbm:tab:phase}, $9.468214$,
$0.894111911$ ($m=10$), $0.815474076$ ($m=20$), $0.775811048$, $2.397371989$
($m=40$), $133.794519$, $0.755843446$, $2.320803680$ ($m=80$). Minimizing the
same proxy over real $n$ in each phase gives $5.929$, $7.319$, $8.734$,
$10.099$ for $m=10,20,40,80$, against $2\rho[\log m+\log\log m+1-\log\rho]=6.243$,
$7.568$, $8.818$, $10.017$: the differences $-0.31$, $-0.25$, $-0.08$, $+0.08$
are consistent with the $o(1)$ of~\eqref{fbm:eq:minimum}, which converges
slowly. At the peak and trough points of the same phases the strict/weak
ratio of the residue sums, times $m/\rho$, is $0.897$, $0.949$, $0.975$,
$0.987$ (peaks) and $0.857$, $0.920$, $0.957$, $0.978$ (troughs), consistent
with the $1+O(1/m)$ of~\eqref{fbm:eq:strictphase}; along the exact sequence
($n\le425$, phases $m\le3$) the same product is still only $0.48$--$0.68$.
\end{writenote}
\begin{writenote}[Correction at the independent check, 7 October 2026]
The range $0.48$--$0.68$ is that of the write's samples $n=50,100,200,300,425$
($0.479\ldots$ to $0.677\ldots$). Over every $n$ with $50\le n\le425$ (phases
$m=2,3$), $(u_n/a_n)\,m/\rho$ runs from $0.4697\ldots$ at $n=124$, the last
index of phase~2, to $0.7041\ldots$ at $n=125$, the first of phase~3: the
factor $m$ jumps at the phase boundary. The conclusion is unchanged.
\end{writenote}

\subsection{Reproduction commands and deterministic packaging}
From the extracted archive, run
\begin{center}
\texttt{python3 reproduce.py -{}-out /tmp/report208-normal}\par
\texttt{python3 -O reproduce.py -{}-out /tmp/report208-optimized}
\end{center}
Each command executes both finite checkers, compares generated data and tables with the shipped versions, rebuilds the PDF to byte stability with \texttt{pdflatex}, and constructs a deterministic ZIP from the actual rebuilt members. The driver omits only the independent checker's deliberately recorded optimization flag when creating its canonical public data file; the raw normal and optimized independent logs are not claimed byte-identical. With the pinned dependencies and \LaTeX\ environment used for this report, the actual final ZIP members and whole ZIP bytes were verified identical across the two modes. Byte identity across different \LaTeX\ distributions is not promised. SHA256 hashes accompany the final deliverables outside the ZIP.

The numerical tests are finite validation. The proofs establish the infinite asymptotic conclusions; neither replaces the other. No claim of certified finite-input inversion, numerical root intervals, or exhaustive literature priority follows from successful reproduction.
\begin{writenote}[The shipped files, 7 October 2026]
In ProveIt the PDF is not shipped; the programs are \file{code/*.py}, their
outputs, the requirements and the generated tables are under \file{data/}.
The programs write their outputs next to themselves and the driver expects the
delivered flat layout (with \file{tables/}), so it runs only in a re-extracted
archive; the report README gives commands that rerun the two checkers on a
scratch copy.
\end{writenote}

\section{Sources and scope of the contribution}\label{fbm:sec:sources}
The definitions, EGFs, and current displayed terms in A386374 and A386375 were inspected directly on 4 October 2026 \cite{weak,strict}. Both entries showed a b-file link but no asymptotic or external research reference. A previously considered cross-reference, A308876, was also inspected: it concerns the EGF $e^z(1-z)/(1-2z)$ and an asymptotic $n!2^{n-1}e^{1/2}$ attributed there to V\'aclav Kot\v{e}\v{s}ovec in 2019, rather than a maximum-occupancy theorem \cite{unrelated}.

The most directly relevant prior is Gourdon's 1996 thesis, Chapter VI, Section 4, Theorem 7 and Example 8 on printed page 195 \cite{gourdonthesis}. Example 8 explicitly treats random surjections with a variable number of image elements, the unrestricted ordered-partition model here. It gives a moving-pole factorial-tail maximum-size distribution, the scale $\log n/\log\log n$, and bounded variance. Thus the factorial cutoff scale, pole perturbation, and leading maximum-size phenomenon are established prior. The shrinking transition width associated with the factorial tail is also a consequence of that established law, rather than a new general extreme-value mechanism.

The actual 1995 proceedings text was also inspected, including the assumptions and Section 5, Theorem 4, printed pages 258--259 \cite{gourdon95}. Its displayed algebraic-logarithmic supercritical theorem assumes a finite singularity of the inner component generating function; $e^z-1$ does not meet that assumption directly. The thesis example handles the entire-function variant. The proceedings also advertises more complete asymptotic expansions. Accordingly, neither pole perturbation nor all-order singularity expansions are presented as a new general method. The 1998 journal version was identified but its full text was not obtained; its contents are not assumed identical to the inspected proceedings \cite{gourdon98}.

Flajolet and Sedgewick, Section V.2 and Notes V.2--V.3, printed page 300, explain renewal conditioning, supercritical sequences, largest components, and truncation-pole displacement, citing Gourdon \cite{ac}. This supplies additional classical context. The displayed smooth coefficient hypothesis there is geometric times a power; it does not itself state the factorial-tail tie expectation and envelopes proved here.

Daly's full manuscript gives logarithmic and Poisson total-variation approximations for maxima multiplicities in a fixed number of independent observations \cite{daly}. Its introduction, equations (1)--(3), and Theorems 1 and 3 were inspected. In the present setting block sizes are constrained by their total and $K$ is random, so \eqref{fbm:eq:expectations} is not an immediate iid substitution. No novelty is claimed for general discrete ties or Poisson approximation. Carayol and Rotondo's full manuscript was inspected in its introduction, fixed-$k$ constrained-Poisson representation, Appendix F maximum bounds, and Appendix G saddle estimates \cite{carayol}. That work gives surjection sampling algorithms and fixed-$k$ profile bounds. The constrained-Poisson representation is prior, rather than a result of this report.

Bender and Gao's paper is particularly relevant because its abstract advertises large parts and multiplicity statistics \cite{bendergao}. Only its publisher abstract and references were obtained; attempted PDF links returned the landing page. Its full text was not read, so theorem-level non-overlap cannot be certified. Olofsson's 1999 article was likewise known through bibliographic and abstract information and Daly's discussion, not a retrieved full original \cite{olofsson}.

The specialized results proved here are the explicit positive-residue coefficient errors, the uniform finite root-free scheme, the first-block tie probability with its two distinct envelope scales and refined phase minimum, the phase-uniform strict/weak comparison, and the shrinking weak inverse enclosure. The inspected passages of Gourdon do not state this collection of first-block enumeration results. This is a bounded comparison with inspected sources, not a first-proof or worldwide-novelty claim. The absence of an asymptotic from a current OEIS entry likewise does not establish originality.

\begin{remark}[{[write]} The OEIS entries, 7~October 2026]\label{fbm:rem:oeis}
The write read the entries in the internal format on 7~October 2026.
\emph{A386374}, revision \#16 (27 July 2025), by John Tyler Rascoe (19 July
2025): ``Number of words of length n over an infinite alphabet such that the
letters cover an initial interval and the letter 1 occurs at least as many
times as any other letter.'', with the e.g.f.\ \eqref{fbm:eq:weakEGF} (as
$\sum_{i\ge0}x^i/(i!\,(1-\sum_{j=1}^ix^j/j!))$), 22 displayed terms, a Maple
program by Alois P.\ Heinz (19 July 2025), a PARI program, and a b-file by
Heinz for $n=0,\ldots,425$. \emph{A386375}, revision \#12 (27 July 2025), by
Rascoe (19 July 2025), the strict version, with the e.g.f.\
\eqref{fbm:eq:strictEGF} in the same form, 23 displayed terms and a b-file by
Heinz for $n=0,\ldots,425$. Neither entry has an asymptotic formula, a
conjecture or a literature reference, as the source states. \emph{A308876},
revision \#15 (18 September 2019), by Ilya Gutkovskiy (29 June 2019),
``Expansion of e.g.f.\ exp(x)*(1 - x)/(1 - 2*x).'', with the formula
$a(n)\sim n!\,2^{n-1}e^{1/2}$ signed by V\'aclav Kote\v{s}ovec, 29 June 2019,
as the source states; it is unrelated to block maxima. \emph{Checks.} All 426
terms of each b-file equal the write's evaluation
of~\eqref{fbm:eq:exactrec}. No OEIS edit is part of this work.
\end{remark}

\subsection{Further calculations}
Three useful next calculations are: explicit constants and onset in Theorem \ref{fbm:thm:expansion}, with interval roots and factorial-tail bounds for finite-threshold certification; higher uniform phase corrections beyond the $o(1)$ phase-minimum remainder; and a separately proved strict-sequence interpolation and integer inverse, respecting its smaller normalized lower bound and finite prefix. These are extensions suggested by this report, not assertions that the questions are unresolved throughout the literature.
\begin{writenote}[Added 7 October 2026, under Vladimir's standing rule of 4 October 2026]
These three calculations are the source's open questions; the write adds from
the non-claims a theorem-level comparison with Bender--Gao's full text (only
its abstract was read) and interval-certified values for the diagnostic
tables. No claim of the source was found false, and the OEIS entries contain
no conjecture.
\end{writenote}
\begin{writenote}[Independent check of the write, 7 October 2026]
An independent adversarial check of the write read A386374 (\#16), A386375
(\#12) and A308876 (\#15) again and recomputed every number the write added,
with its own code.
\begin{itemize}
\item The counts equal both b-files (426 terms each). They also equal the
coefficients of the two OEIS e.g.f.s, expanded as exact rational series to
degree $60$, a route that does not use~\eqref{fbm:eq:exactrec}. A brute
force over all initial-alphabet words with $n\le7$ gives $a_n$, $u_n$ and
$P_n=\E[J/K]$ exactly, and $a_{n+1}>a_n$ holds for $1\le n<425$.
\item The numerical constants of the proofs and the two rational checks
hold. A SymPy reversion of its own gives the coefficients
\eqref{fbm:eq:ellcoeff} and~\eqref{fbm:eq:Q12}. The factorization
$c_m(\rho/r_m)^n=T_m(n)G_{m,nD_m}(1/n)$ is an exact identity, since $d_m(D_m)=r_m-\rho$.
\item Tables~\ref{fbm:tab:diagnostic} and~\ref{fbm:tab:phase} agree in every
printed digit, with $D_m$ summed as a tail and the roots taken at 100 and 320
digits; every truncation in the write's note is correct. The phase minima
($5.929$, $7.319$, $8.734$, $10.099$) and the strict/weak ratios at the
peaks and troughs agree too.
\end{itemize}
It also confirmed the provenance figures (470,406 bytes, 13 files, 551 lines,
15 pages), the byte identity of the ten staged files, the identity of the three inlined tables with
the shipped table files, Remark~\ref{fbm:rem:transseries}, the label numbering
(68 delivered labels unchanged, 71 in all) and the README listing. Two
statements of the write needed precision, given in dated notes: the Bell-pole
deviation bound (after the provenance note) and the range of the strict/weak
product along the exact sequence (after the tables' note). No claim was found
wrong in substance.
\end{writenote}

\begin{thebibliography}{99}
\small\setlength{\itemsep}{3pt}
\bibitem{weak} OEIS Foundation Inc. \emph{A386374}. Current entry inspected 4 October 2026. \url{https://oeis.org/A386374}.
\bibitem{strict} OEIS Foundation Inc. \emph{A386375}. Current entry inspected 4 October 2026. \url{https://oeis.org/A386375}.
\bibitem{unrelated} OEIS Foundation Inc. \emph{A308876}. Current entry inspected 4 October 2026. \url{https://oeis.org/A308876}.
\bibitem{gourdonthesis} X. Gourdon. \emph{Combinatoire, Algorithmique et G\'eom\'etrie des Polyn\^omes}. Doctoral thesis, \'{E}cole Polytechnique, 1996. Chapter VI, Section 4 and Example 8, printed p.\ 195. \href{https://www.mit.edu/~sfinch/Xavier_Gourdon_-_Combinatoire,_algorithmique_et_geometrie_des_polynomes_-_Ecole_Polytechnique_thesis.pdf}{Full thesis PDF}.
\bibitem{gourdon95} X. Gourdon. Largest component in random combinatorial structures. \emph{FPSAC 1995 proceedings}, 1995, 12-page proceedings version; Section 5, Theorem 4, printed pp.\ 258--259. \url{https://fpsac-archive.github.io/FPSAC95/ARTICLES/24.GOURDON.pdf}.
\bibitem{gourdon98} X. Gourdon. Largest component in random combinatorial structures. \emph{Discrete Mathematics} 180 (1998), 185--209. \url{https://doi.org/10.1016/S0012-365X(97)00115-5}. Journal full text not obtained for this comparison.
\bibitem{ac} P. Flajolet and R. Sedgewick. \emph{Analytic Combinatorics}. Cambridge University Press, 2009. Section V.2 and Notes V.2--V.3. \url{https://algo.inria.fr/flajolet/Publications/book.pdf}.
\bibitem{daly} F. Daly. \emph{Approximations for the number of maxima and near-maxima in independent data}. arXiv:2505.06088v1. Inspected HTML displayed manuscript date 24 August 2026. \url{https://arxiv.org/html/2505.06088v1}.
\bibitem{carayol} A. Carayol and P. Rotondo. \emph{Efficient Uniform Sampling of Surjections via their Profiles}. arXiv:2605.24068v1, 22 May 2026. \url{https://arxiv.org/html/2605.24068v1}.
\bibitem{bendergao} E. A. Bender and Z. Gao. Part Sizes of Smooth Supercritical Compositional Structures. \emph{Combinatorics, Probability and Computing} 23 (2014), 686--716. \url{https://doi.org/10.1017/S0963548314000315}. Publisher abstract and references inspected; full text not obtained.
\bibitem{olofsson} P. Olofsson. A Poisson approximation with applications to the number of maxima in a discrete sample. 1999. Bibliographic and abstract lead discussed by Daly \cite{daly}; original full article not obtained.
\bibitem{TSvol}{\raggedright ProveIt, \emph{Transseries and inversion}, \file{Analysis/Transseries/docs/series-and-transseries/Transseries_And_Inversion/transseries_and_inversion.tex}; Part~0 (\file{p0:def:model}, \file{p0:thm:lambert-core}, \file{p0:def:three-inverses}, \file{p0:thm:staircase}) and the chapter on the Lambert template (\file{plt:thm:lw-template}). [Added in the write.]\par}
\end{thebibliography}
\end{document}
